\section{What migrant women describe}
\label{sec:migrant-findings}
The accounts distinguish a payment that someone can make from a service she can use on her own terms. A wallet can be useful for daily purchases or receiving a customer's payment without giving a woman control over the account, ensuring that money reaches a relative abroad, or giving her a workable way to seek redress. We develop these distinctions from Angélica Cardozo Cadavid's thematic synthesis and accompanying memo on interviews and focus groups. We use the range of arrangements that the synthesis describes; we do not count mentions, attribute a view to an entire group, or reconstruct which individuals contributed to each theme.

\subsection{A domestic payment can be only the first leg}
\label{sec:migrant-journeys}

The migrant synthesis (Sections 3.1--3.6) describes women buying goods, collecting customer payments and using transfers, wallets or QR codes in their businesses. A local payment can complete a purchase in Colombia. In the synthesis's accounts of sending toward Venezuela, a local payment can instead be the first leg. The written accounts describe checking an exchange rate with a contact, paying into a Colombian account and separately checking delivery abroad. These steps illustrate a route described across the accounts, not one timed or independently verified transaction. The source does not give a common clock for the time between the local payment and usable delivery.

The migrant memo on interviews and focus groups (first thematic synthesis, Section 5) also describes inward transfers that women experienced as faster or more direct. These are comparisons in the written analysis, not measured travel times or a matched set of inward and outward transfers. That variation is important. An arrangement that helps someone send money toward Venezuela need not be the arrangement someone else used to receive money from abroad. The survey's disposition question asks how incoming proceeds generally end up, not whether the respondent completed an outbound transfer. The accounts describe how domestic payment use and a need for help abroad can coexist; they do not rank inbound and outbound services or explain the numerical difference between the survey's two outcomes.

\subsection{Intermediaries provide a service and can create a dependency}

Intermediaries could arrange exchange and delivery that the sender's own wallet did not provide. Women described using recommendations, returning to contacts they already knew and comparing the terms that different contacts offered. A relationship can be valuable in its own right: a known contact can explain a step or confirm delivery. This is why it would be misleading to describe every use of an intermediary as a failure of digital access. The same arrangement can be hard to replace, particularly when the sender does not know who handles later steps or what another contact would charge. The accounts support both interpretations; they do not tell us how many women chose an intermediary over an equally workable direct route.

The extra work matters even when the first payment is fast. The synthesis describes checking rates and which account to send to, coordinating across people and finding out whether the recipient got usable funds. A stated fee is not the same as the exchange-rate cost, and the time for the local payment omits the wait elsewhere in the chain. The reports of layered charges and uncertainty illustrate what a total-price or time-to-delivery measure would need to capture, not how much every sender loses. Women also described using payments to meet household needs and keep small businesses running. A delay can matter for an obligation even if the payment eventually arrives; the study does not quantify that welfare cost.

\subsection{Entry, assistance and control are separate}
\label{sec:migrant-control}

The written migrant synthesis (Sections 3.8 and 4.4--4.5) distinguishes two barriers to entry: not holding a document a provider requested and holding one that a provider or app reportedly did not accept. Helping someone obtain a document may address the first; it need not address the second. A reported line or proof of address in one's own name is therefore informative about onboarding conditions without verifying that a specific product was available. The institutional interviews help explain what verification problems could look like, but cannot tell us which rule determined the outcome for any woman in the study (Section~\ref{sec:expert-findings}).

The synthesis also describes payments through accounts held by a relative or partner (Section 4.5). This is a second, domestic example of how a payment can work with assistance. The synthesis does not establish a particular participant, transaction date or recall period for the example. Using the account can make a purchase possible without giving the user the same control as its holder. We distinguish initiating a payment, having control of the account or its credentials, and deciding how household money is spent. The evidence does not establish that all shared arrangements are involuntary or that an individually held account guarantees discretion over household funds. This distinction matters for interpreting agency: access to a resource and the ability to act on a decision are different \citep{kabeer1999}.

The accounts include practical knowledge gained through transacting and shared with others. That makes learning by doing a plausible reading of the survey association between reported ability and domestic use. Prior skills could also make it easier to start using a service, and assistance can make a transaction possible before someone can complete it unaided. Neither the synthesis nor the survey separates those explanations across respondents. We therefore read the ability score as a report of capability and confidence, not a skills test or an estimate of what training would do.

\subsection{Taking precautions does not secure a remedy}
\label{sec:migrant-recourse}

The migrant synthesis (Section 3.7) and the first thematic memo (Sections 7--8) describe suspicious links and false payment confirmations. Women reported checking requests or receipts, using known contacts and asking relatives to verify arrival. These steps can answer an immediate question---for example, whether the family received money---without identifying the provider or intermediary who should act if something went wrong. Trust in a known person can help someone attempt the transfer. It does not assure a prompt response or a refund.

This matters for how we interpret both sources. The accounts concern reported experiences and protective practices, not independently verified fraud or recovery. In the survey, reporting a problem opens the question about complaint channels. The first-response and resolution questions concern respondents who complained; the latter does not require a reported first response. A response does not establish restitution. The retained focus-group guide also includes a hypothetical Bre-B transfer. A response to that exercise describes expectations or concerns, not a transaction carried out on Bre-B. These distinctions keep users' experiences visible without attributing unverified failures or successes to a named rail.

\subsection{The qualitative positions describe constraints, not types}
\label{sec:migrant-positions}

The earlier synthesis organized its interpretation around four overlapping positions: barriers to formal entry; reliance on intermediaries; practical digital capability constrained by the services on offer; and more independent use that still faces cross-border constraints. Safety concerns cut across them. A woman can transact independently for one purpose and rely on help for another. The positions therefore help ask where autonomy ends, not how many discrete kinds of users exist. The four observed survey configurations in Figure~\ref{fig:onboarding-conditions} pair two reported conditions; neither those configurations nor any fitted class label is assigned to an interviewee. The joint display uses that distinction to connect the findings without merging the samples.
